% !TeX root = main.tex
\cleardoublepage
\phantomsection
\chapter*{Conclusion générale}
\addcontentsline{toc}{chapter}{Conclusion générale}

Ce projet de fin d’études a porté sur l’évolution de la plateforme \textit{HireCue}, une solution SaaS dédiée au recrutement intelligent. L’objectif principal était d’améliorer la plateforme existante en renforçant ses capacités d’analyse, d’automatisation et d’aide à la décision.

Les travaux réalisés ont permis d’enrichir la plateforme à travers plusieurs modules complémentaires. Les fonctionnalités basées sur l’intelligence artificielle ont contribué à rendre le processus de recrutement plus explicable, plus personnalisé et plus efficace pour les différents utilisateurs de la plateforme.

L’automatisation a constitué un autre axe important du projet. La mise en place de workflows a permis de réduire certaines tâches répétitives, de faciliter l’intégration de données externes et d’améliorer la traçabilité des opérations. Ces mécanismes participent à une meilleure efficacité opérationnelle et à une gestion plus fluide des processus de recrutement.

Par ailleurs, l’intégration d’une couche décisionnelle a permis de valoriser les données générées par la plateforme. Les tableaux de bord développés offrent une vision synthétique de l’activité, facilitent le suivi des indicateurs clés et soutiennent la prise de décision à différents niveaux.

Ce projet a également constitué une opportunité de mettre en pratique des compétences variées liées au développement logiciel, à l’intégration de services, à l’intelligence artificielle, à l’automatisation des processus et à l’analyse de données. Il a permis de mieux comprendre les enjeux techniques et organisationnels associés à la conception d’une plateforme de recrutement moderne.

Enfin, plusieurs perspectives restent envisageables afin de renforcer davantage la plateforme, notamment l’amélioration des fonctionnalités d’intelligence artificielle, l’évolution des mécanismes d’automatisation et l’enrichissement de la couche décisionnelle. Ces évolutions permettront de poursuivre le développement de HireCue tout en conservant l’objectif principal du projet : mettre la technologie au service d’un recrutement plus efficace, plus transparent et plus centré sur l’humain.